% --- Tabel Deskripsi UC-04 ---
\begin{table}[H]
    \centering
    \caption{Deskripsi Use Case Melihat Umpan Balik Rekan (UC-04)}
    \label{tbl:desc_uc04}
    \begin{tabular}{|p{0.25\textwidth}|p{0.70\textwidth}|}
        \hline
        \textbf{Nama Use Case} & \textbf{Melihat Umpan Balik Rekan} \\
        \hline
        \textbf{Aktor} & Siswa \\
        \hline
        \textbf{FR Terkait} & FR-03 \\
        \hline
        \textbf{Deskripsi} & Siswa meninjau hasil evaluasi dan komentar yang diberikan oleh rekan sejawat terhadap tugas yang telah dikerjakan sebelumnya. \\
        \hline
        \textbf{Kondisi sebelum} & Proses \textit{peer assessment} oleh rekan telah selesai; Notifikasi tersedia. \\
        \hline
        \textbf{Kondisi sesudah} & Status notifikasi umpan balik menjadi "terbaca". \\
        \hline
        \textbf{Main Flow} & \begin{enumerate}[noitemsep,topsep=0pt,leftmargin=*]
            \item Siswa memilih notifikasi hasil asesmen.
            \item Sistem menampilkan rincian skor dan komentar dari rekan.
            \item Siswa mempelajari masukan yang diberikan.
            \item Siswa menutup tampilan umpan balik.
        \end{enumerate} \\
        \hline
    \end{tabular}
\end{table}